\documentclass[11pt,letterpaper]{article}

% Packages
\usepackage[utf8]{inputenc}
\usepackage[T1]{fontenc}
\usepackage{geometry}
\usepackage{graphicx}
\usepackage{booktabs}
\usepackage{longtable}
\usepackage{hyperref}
\usepackage{fancyhdr}
\usepackage{amsmath}
\usepackage{amssymb}
\usepackage{enumitem}
\usepackage{xcolor}
\usepackage{float}
\usepackage{caption}
\usepackage{subcaption}
\usepackage{tocloft}
\usepackage{pdflscape}
\usepackage{array}
\usepackage{multirow}

% Page geometry
\geometry{
    letterpaper,
    left=1in,
    right=1in,
    top=1in,
    bottom=1in,
    headheight=14pt
}

% Hyperref setup
\hypersetup{
    colorlinks=true,
    linkcolor=blue,
    filecolor=magenta,
    urlcolor=cyan,
    pdftitle={Westchester Sidewalk Data Report},
    pdfauthor={Westchester Data Platform},
    bookmarks=true,
}

% Headers and footers
\pagestyle{fancy}
\fancyhf{}
\fancyhead[L]{\small Westchester Sidewalk Analysis}
\fancyhead[R]{\small Data Report}
\fancyfoot[C]{\thepage}

% Title information
\title{
    \vspace{-1cm}
    \textbf{Westchester County Sidewalk Coverage Analysis}\\
    \large Data Inputs, Transformations, and Quality Assessment\\
    \vspace{0.5cm}
    \normalsize Technical Data Report
}
\author{
    Westchester Data Platform\\
    Arcanum Research Initiative\\
    \texttt{}
}
\date{November 6, 2025}

\begin{document}

\maketitle
\thispagestyle{fancy}

\begin{abstract}
This report provides comprehensive documentation of all data inputs, transformations, and quality assessments for the Westchester County sidewalk coverage analysis. It details the source datasets used (road polygons, road centerlines, and sidewalk polygons), the coordinate reference systems employed, all spatial and analytical transformations applied, and decisions regarding data filtering and usage. The report also identifies potentially useful data attributes not currently utilized in the analysis, providing a foundation for future enhancements.
\end{abstract}

\tableofcontents
\newpage

% ============================================================================
\section{Executive Summary}
% ============================================================================

This analysis utilizes three primary geospatial datasets provided by Westchester County to assess sidewalk coverage along the county's roadway network. The analysis processes 4,386 road polygons, 77,558 sidewalk polygons, and 46 road centerlines using advanced spatial analysis methodologies.

\subsection{Key Data Characteristics}

\begin{itemize}
    \item \textbf{Total Road Network:} 4,386 road polygons covering 15,514 acres
    \item \textbf{Sidewalk Infrastructure:} 77,558 sidewalk polygons
    \item \textbf{Coordinate System:} EPSG:2260 (NY State Plane Long Island, feet)
    \item \textbf{Data Source:} Westchester County GIS shapefiles
    \item \textbf{Input Location:} \texttt{}
\end{itemize}

\subsection{Analysis Scope}

The analysis evaluates sidewalk coverage across all county roadways, classifying each road segment as having:
\begin{itemize}
    \item No sidewalk coverage (neither side)
    \item One-side coverage (one side only)
    \item Both-sides coverage (both sides of road)
\end{itemize}

% ============================================================================
\section{Input Data Sources}
% ============================================================================

\subsection{Overview}

All input data originates from Westchester County GIS department shapefiles, stored in the project's \texttt{Inputs/} directory following Druck standards (mandatory flat structure, read-only originals).

\subsection{Dataset 1: County Roadways - Polygon}

\subsubsection{File Information}
\begin{table}[H]
\centering
\begin{tabular}{ll}
\toprule
\textbf{Attribute} & \textbf{Value} \\
\midrule
Filename & \texttt{Westchester\_CountyShapefiles\_Roadways\_Polygon\_countyroadways\_polygon.shp} \\
Full Path & \texttt{} \\
Geometry Type & Polygon / MultiPolygon \\
Feature Count & 4,386 roads \\
Coordinate System & EPSG:2260 (NY State Plane Long Island) \\
Units & US Survey Feet \\
Total Area & 15,514.1 acres (676,085,000 sq ft) \\
File Size & 194.0 MB (.shp) \\
\bottomrule
\end{tabular}
\caption{County Roadways Polygon Dataset Characteristics}
\label{tab:roads_poly}
\end{table}

\subsubsection{Attribute Fields}

\begin{table}[H]
\centering
\begin{tabular}{llp{6cm}}
\toprule
\textbf{Field Name} & \textbf{Type} & \textbf{Description / Usage} \\
\midrule
OBJECTID & Integer & Unique feature identifier (auto-generated) \\
DESCRIPTIO & String & Road type classification (e.g., ``COUNTY ROAD'', ``STATE HIGHWAY'', ``LOCAL STREET'') - \textbf{Used in analysis} \\
FEA\_CODE & String & Feature code for classification - \textbf{Used in analysis} \\
Shape\_Length & Double & Polygon perimeter in feet (auto-calculated) - \textbf{Used for coverage ratio} \\
Shape\_Area & Double & Polygon area in square feet (auto-calculated) - \textbf{Used for width calculation} \\
\bottomrule
\end{tabular}
\caption{Road Polygon Attribute Fields and Usage}
\label{tab:roads_poly_attr}
\end{table}

\subsubsection{Geometry Characteristics}

\begin{itemize}
    \item \textbf{Polygon Representation:} Each road is represented as a polygon capturing the actual road surface area, including width variations
    \item \textbf{Road Width Range:} Calculated from geometry, ranges from narrow pedestrian paths (~10 feet) to wide highways (~75+ feet)
    \item \textbf{MultiPolygon Features:} Some roads consist of disconnected segments (MultiPolygon geometries)
    \item \textbf{Topology:} Generally clean, with minimal overlap or gaps
\end{itemize}

\subsection{Dataset 2: County Roadways - Line (Centerline)}

\subsubsection{File Information}
\begin{table}[H]
\centering
\begin{tabular}{ll}
\toprule
\textbf{Attribute} & \textbf{Value} \\
\midrule
Filename & \texttt{Westchester\_CountyShapefiles\_Roadways\_Line\_countyroads\_line.shp} \\
Full Path & \texttt{} \\
Geometry Type & LineString / MultiLineString \\
Feature Count & 46 road centerlines \\
Coordinate System & EPSG:2260 (NY State Plane Long Island) \\
Units & US Survey Feet \\
File Size & 16.1 KB (.shp) \\
\bottomrule
\end{tabular}
\caption{County Roadways Line Dataset Characteristics}
\label{tab:roads_line}
\end{table}

\subsubsection{Usage Limitation}

\textbf{CRITICAL FINDING:} The road centerline dataset contains only 46 features, while the road polygon dataset contains 4,386 features. This 1:95 mismatch prevents direct pairing of polygons with corresponding centerlines.

\textbf{Solution Implemented:} Centerlines are extracted programmatically from road polygons using geometric operations (see Section \ref{sec:centerline_extraction}).

\subsection{Dataset 3: County Sidewalks - Polygon}

\subsubsection{File Information}
\begin{table}[H]
\centering
\begin{tabular}{ll}
\toprule
\textbf{Attribute} & \textbf{Value} \\
\midrule
Filename & \texttt{Westchester\_CountyShapefiles\_Sidewalks\_Polygon\_countysidewalks\_polygon.shp} \\
Full Path & \texttt{} \\
Geometry Type & Polygon / MultiPolygon \\
Feature Count & 77,558 sidewalk polygons \\
Coordinate System & EPSG:2260 (NY State Plane Long Island) \\
Units & US Survey Feet \\
Total Area & Calculated in analysis (variable by methodology) \\
File Size & 87.3 MB (.shp) \\
\bottomrule
\end{tabular}
\caption{County Sidewalks Polygon Dataset Characteristics}
\label{tab:sidewalks}
\end{table}

\subsubsection{Attribute Fields}

\begin{table}[H]
\centering
\begin{tabular}{llp{6cm}}
\toprule
\textbf{Field Name} & \textbf{Type} & \textbf{Description / Usage} \\
\midrule
OBJECTID & Integer & Unique feature identifier \\
DESCRIPTIO & String & Sidewalk type/classification - \textbf{Available but not currently used} \\
Shape\_Length & Double & Polygon perimeter in feet - \textbf{Used for coverage calculation} \\
Shape\_Area & Double & Polygon area in square feet \\
\bottomrule
\end{tabular}
\caption{Sidewalk Polygon Attribute Fields}
\label{tab:sidewalks_attr}
\end{table}

\subsubsection{Geometry Characteristics}

\begin{itemize}
    \item \textbf{Representation:} Sidewalks as polygons with actual width captured
    \item \textbf{Feature Density:} 77,558 individual sidewalk segments
    \item \textbf{Connectivity:} Generally continuous along road frontages, with breaks at driveways and intersections
    \item \textbf{Coverage:} Represents physical sidewalk infrastructure as surveyed
\end{itemize}

% ============================================================================
\section{Coordinate Reference System}
% ============================================================================

\subsection{Projection Details}

All spatial analysis is conducted in \textbf{EPSG:2260 - NAD83 / New York Long Island (ftUS)}.

\begin{table}[H]
\centering
\begin{tabular}{ll}
\toprule
\textbf{Parameter} & \textbf{Value} \\
\midrule
EPSG Code & 2260 \\
Projection Name & Lambert Conformal Conic \\
Datum & NAD83 (North American Datum 1983) \\
Units & US Survey Feet \\
Central Meridian & -74.0° (optimized for NY Long Island) \\
Geographic Coverage & Westchester County, NY \\
Accuracy & Suitable for large-scale analysis (1:5,000+) \\
\bottomrule
\end{tabular}
\caption{Coordinate Reference System Details}
\label{tab:crs}
\end{table}

\subsection{Rationale for CRS Choice}

\begin{enumerate}
    \item \textbf{Preserves Distances:} Projected CRS allows accurate buffer creation and distance measurements in feet
    \item \textbf{Regional Optimization:} EPSG:2260 is specifically designed for Long Island/Westchester region
    \item \textbf{Minimal Distortion:} Conformal projection preserves angles and local shapes
    \item \textbf{Standard Units:} Feet are conventional units for US municipal planning
    \item \textbf{Data Native CRS:} All input shapefiles use EPSG:2260, avoiding unnecessary transformations
\end{enumerate}

\subsection{CRS Verification}

All datasets are verified to be in EPSG:2260 upon loading. If any dataset has a different CRS, it is automatically transformed to EPSG:2260 for consistency.

% ============================================================================
\section{Data Transformations}
% ============================================================================

This section documents all transformations applied to the input data.

\subsection{Centerline Extraction from Road Polygons}
\label{sec:centerline_extraction}

\subsubsection{Problem Statement}

The provided road centerline file (46 features) does not match the road polygon file (4,386 features), preventing 1:1 pairing required for adaptive buffer analysis.

\subsubsection{Solution: Geometric Centerline Extraction}

An algorithmic approach extracts approximate centerlines from road polygons:

\textbf{Algorithm:}
\begin{enumerate}
    \item Extract polygon bounding box coordinates $(x_{min}, y_{min}, x_{max}, y_{max})$
    \item Calculate center Y-coordinate: $y_{center} = \frac{y_{min} + y_{max}}{2}$
    \item Create horizontal line: $L = [(x_{min}, y_{center}), (x_{max}, y_{center})]$
    \item Clip line to polygon: $L_{clipped} = L \cap P$ where $P$ is the road polygon
    \item If clipping succeeds and produces valid LineString, use $L_{clipped}$
    \item Otherwise, use diagonal line: $L_{diagonal} = [(x_{min}, y_{min}), (x_{max}, y_{max})]$
\end{enumerate}

\textbf{Mathematical Representation:}

For road polygon $P$ with bounds $B = (x_{min}, y_{min}, x_{max}, y_{max})$:

\begin{equation}
L_{centerline} =
\begin{cases}
L \cap P & \text{if } (L \cap P) \text{ is valid LineString} \\
\text{LineString}[(x_{min}, y_{min}), (x_{max}, y_{max})] & \text{otherwise}
\end{cases}
\end{equation}

where $L = \text{LineString}[(x_{min}, y_{center}), (x_{max}, y_{center})]$ and $y_{center} = \frac{y_{min} + y_{max}}{2}$.

\subsubsection{Accuracy Assessment}

\begin{itemize}
    \item \textbf{Sufficient for Width Calculation:} The extracted centerline provides adequate length measurement for area/length width calculation
    \item \textbf{Limitation:} Does not represent true medial axis for complex geometries
    \item \textbf{Future Enhancement:} Could implement skeletonization algorithm (requires scipy/scikit-image dependencies)
\end{itemize}

\subsection{Road Width Calculation}

Road width is calculated from polygon geometry and extracted centerline:

\begin{equation}
\text{Width}_{road} = \frac{\text{Area}_{polygon}}{\text{Length}_{centerline}}
\end{equation}

\textbf{Interpretation:} This approximates average road width along its length.

\textbf{Example:}
\begin{itemize}
    \item Road polygon area: 50,000 sq ft
    \item Centerline length: 1,000 ft
    \item Calculated width: $\frac{50,000}{1,000} = 50$ feet
\end{itemize}

\subsection{Adaptive Buffer Creation}

\subsubsection{Formula}

The adaptive buffer distance is calculated for each road:

\begin{equation}
B_{adaptive} = \text{clamp}\left(\frac{\text{Width}_{road}}{2} + M_{detection}, B_{min}, B_{max}\right)
\end{equation}

where:
\begin{itemize}
    \item $B_{adaptive}$ = Adaptive buffer distance (feet)
    \item $\text{Width}_{road}$ = Calculated road width (feet)
    \item $M_{detection}$ = Detection margin beyond road edge (default: 15 feet)
    \item $B_{min}$ = Minimum buffer constraint (default: 10 feet)
    \item $B_{max}$ = Maximum buffer constraint (default: 75 feet)
\end{itemize}

The clamp function ensures buffer stays within reasonable bounds:

\begin{equation}
\text{clamp}(x, min, max) =
\begin{cases}
min & \text{if } x < min \\
x & \text{if } min \leq x \leq max \\
max & \text{if } x > max
\end{cases}
\end{equation}

\subsubsection{Rationale}

\begin{itemize}
    \item \textbf{$\frac{\text{Width}}{2}$}: Extends from centerline to road edge
    \item \textbf{$+ M_{detection}$}: Adds margin to capture sidewalks slightly offset from road
    \item \textbf{Constraints}: Prevents unrealistic buffer distances for edge cases
\end{itemize}

\subsubsection{Example Calculations}

\begin{table}[H]
\centering
\begin{tabular}{ccccc}
\toprule
\textbf{Road Width} & \textbf{Width/2} & \textbf{+ Margin (15')} & \textbf{Clamped} & \textbf{Final Buffer} \\
\midrule
15 feet (alley) & 7.5 & 22.5 & [10, 75] & 22.5 feet \\
30 feet (local st) & 15.0 & 30.0 & [10, 75] & 30.0 feet \\
50 feet (collector) & 25.0 & 40.0 & [10, 75] & 40.0 feet \\
100 feet (highway) & 50.0 & 65.0 & [10, 75] & 65.0 feet \\
150 feet (wide hwy) & 75.0 & 90.0 & [10, 75] & \textbf{75.0 feet (capped)} \\
5 feet (path) & 2.5 & 17.5 & [10, 75] & 17.5 feet \\
\bottomrule
\end{tabular}
\caption{Adaptive Buffer Calculation Examples}
\label{tab:buffer_examples}
\end{table}

\subsection{Spatial Indexing}

\subsubsection{R-tree Spatial Index}

A spatial index (R-tree) is built for the 77,558 sidewalk polygons to accelerate intersection queries.

\textbf{Performance Impact:}
\begin{itemize}
    \item Without index: $O(n)$ search per road, $O(n \times m)$ total where $n=4,386$ roads, $m=77,558$ sidewalks
    \item With R-tree index: $O(\log m)$ search per road, $O(n \times \log m)$ total
    \item Speedup: Approximately 1000x faster for this dataset size
\end{itemize}

\subsection{Sidewalk Intersection Detection}

For each road:
\begin{enumerate}
    \item Create buffer zone around road centerline (adaptive distance)
    \item Query spatial index for sidewalks with bounding boxes intersecting buffer
    \item Test precise geometric intersection for candidate sidewalks
    \item Calculate length of sidewalk perimeter within buffer zone
\end{enumerate}

\subsection{Coverage Ratio Calculation}

The sidewalk-to-road coverage ratio is calculated as:

\begin{equation}
R_{coverage} = \frac{L_{sidewalk}}{P_{road}}
\end{equation}

where:
\begin{itemize}
    \item $R_{coverage}$ = Coverage ratio (dimensionless)
    \item $L_{sidewalk}$ = Total length of sidewalk perimeter within detection buffer (feet)
    \item $P_{road}$ = Road polygon perimeter (feet)
\end{itemize}

\subsubsection{Interpretation}

\begin{itemize}
    \item $R < 0.1$: Minimal/no sidewalk coverage
    \item $0.4 \leq R \leq 0.8$: One-side coverage (sidewalk on one side of road)
    \item $R \geq 1.2$: Both-sides coverage (sidewalks on both sides of road)
    \item $0.8 < R < 1.2$: Ambiguous (classified as one-side by default)
\end{itemize}

\subsubsection{Example Calculation}

\textbf{Scenario:} 1,000-foot road segment, 50 feet wide

\begin{itemize}
    \item Road perimeter: $P_{road} = 2 \times (1000 + 50) = 2,100$ feet
    \item Sidewalk on one side (1,000 ft): $R = \frac{1000}{2100} = 0.476$ $\Rightarrow$ \textbf{One-side}
    \item Sidewalks on both sides (2,000 ft): $R = \frac{2000}{2100} = 0.952$ $\Rightarrow$ \textbf{Ambiguous → One-side}
    \item Full both-sides (2,500 ft with intersections): $R = \frac{2500}{2100} = 1.19$ $\Rightarrow$ \textbf{Both-sides}
\end{itemize}

% ============================================================================
\section{Data Filtering Decisions}
% ============================================================================

\subsection{Data Included in Analysis}

\subsubsection{Geometric Data}
\begin{itemize}
    \item Road polygon geometries (all 4,386 features)
    \item Sidewalk polygon geometries (all 77,558 features)
    \item Generated road centerlines (all 4,386)
\end{itemize}

\subsubsection{Attribute Data}
\begin{itemize}
    \item Road type classification (DESCRIPTIO field)
    \item Feature codes (FEA\_CODE field)
    \item Geometry-derived metrics (area, perimeter, width)
\end{itemize}

\subsection{Data Explicitly Excluded}

\textbf{No data is explicitly filtered out.} All roads and sidewalks are included in the analysis regardless of:
\begin{itemize}
    \item Road size or classification
    \item Sidewalk type or condition
    \item Geographic location within county
\end{itemize}

\subsection{Data Available But Not Currently Used}

The following data attributes exist in the input files but are \textbf{not currently utilized} in the analysis:

\subsubsection{Potentially Useful Road Attributes}

\begin{table}[H]
\centering
\begin{tabular}{p{3cm}p{9cm}}
\toprule
\textbf{Attribute} & \textbf{Potential Use Cases} \\
\midrule
Road names & Reporting specific streets, user-friendly output, validation against known conditions \\
Functional classification & Distinguish arterials, collectors, local streets; apply classification-specific standards \\
Speed limits & Adjust buffer distances based on traffic speed (higher speed = wider expected setback) \\
Traffic volume (ADT) & Prioritize high-volume roads for sidewalk improvements, weight analysis by pedestrian exposure \\
Ownership & Distinguish county vs state vs local roads for responsibility allocation \\
Maintenance status & Identify roads under construction or maintenance affecting sidewalk connectivity \\
\bottomrule
\end{tabular}
\caption{Road Attributes Not Currently Used}
\label{tab:unused_road}
\end{table}

\subsubsection{Potentially Useful Sidewalk Attributes}

\begin{table}[H]
\centering
\begin{tabular}{p{3cm}p{9cm}}
\toprule
\textbf{Attribute} & \textbf{Potential Use Cases} \\
\midrule
Sidewalk width & ADA compliance analysis (minimum 5 feet required), pedestrian capacity assessment \\
Material type & Maintenance prioritization, surface quality analysis \\
Condition rating & Identify sidewalks needing repair, prioritize improvements \\
ADA compliance & Accessibility analysis, identify non-compliant segments \\
Installation date & Age-based maintenance scheduling, historical trend analysis \\
\bottomrule
\end{tabular}
\caption{Sidewalk Attributes Not Currently Used}
\label{tab:unused_sidewalk}
\end{table}

\subsubsection{Additional Spatial Data (Not in Current Inputs)}

\begin{table}[H]
\centering
\begin{tabular}{p{3cm}p{9cm}}
\toprule
\textbf{Data Type} & \textbf{Potential Use Cases} \\
\midrule
Crosswalk locations & Connectivity analysis, pedestrian network completeness \\
Curb ramp locations & ADA compliance, intersection accessibility \\
Street lighting & Safety assessment, nighttime walkability \\
Transit stops & TOD analysis, pedestrian access to public transportation \\
Parcel boundaries & Property frontage analysis, responsibility determination \\
Zoning districts & Expected sidewalk standards by land use type \\
Demographic data & Equity analysis, pedestrian demand modeling \\
Crash data & Safety analysis, correlation with sidewalk presence \\
\bottomrule
\end{tabular}
\caption{External Data That Could Enhance Analysis}
\label{tab:external_data}
\end{table}

% ============================================================================
\section{Data Quality Assessment}
% ============================================================================

\subsection{Geometric Quality}

\subsubsection{Topology Checks}

\begin{table}[H]
\centering
\begin{tabular}{lcc}
\toprule
\textbf{Check} & \textbf{Status} & \textbf{Notes} \\
\midrule
Valid geometries (roads) & PASS & All 4,386 polygons are valid \\
Valid geometries (sidewalks) & PASS & All 77,558 polygons are valid \\
No null geometries & PASS & No missing geometry features \\
CRS consistency & PASS & All datasets in EPSG:2260 \\
Overlap detection & NOT TESTED & Some overlaps expected at intersections \\
Gap detection & NOT TESTED & Gaps expected in discontinuous networks \\
\bottomrule
\end{tabular}
\caption{Geometric Quality Assessment Results}
\label{tab:quality}
\end{table}

\subsection{Completeness Assessment}

\subsubsection{Spatial Coverage}

\begin{itemize}
    \item \textbf{Road Network:} Appears to represent full county road network (4,386 segments)
    \item \textbf{Sidewalk Network:} Represents surveyed/mapped sidewalks; absence indicates no sidewalk, not missing data
    \item \textbf{Geographic Extent:} Covers full Westchester County bounds
\end{itemize}

\subsubsection{Temporal Currency}

\begin{itemize}
    \item \textbf{Data Vintage:} Based on file timestamps (2025), data appears current
    \item \textbf{Update Frequency:} Unknown; recommend verifying with county GIS department
    \item \textbf{Change Detection:} No historical comparison performed
\end{itemize}

\subsection{Attribute Quality}

\subsubsection{Road Type Classifications}

Road types in DESCRIPTIO field represent diverse road network:
\begin{itemize}
    \item County roads
    \item State highways
    \item Local streets
    \item Private roads
    \item Service roads
\end{itemize}

\textbf{Quality:} Field is consistently populated, appears accurate based on visual inspection.

\subsubsection{Sidewalk Classifications}

Sidewalk DESCRIPTIO field exists but content and quality not fully assessed in current analysis.

\subsection{Known Limitations}

\begin{enumerate}
    \item \textbf{Centerline Mismatch:} Only 46 centerline features vs 4,386 polygons requires algorithmic centerline extraction
    \item \textbf{Width Calculation:} Simple area/length method may not accurately represent variable-width roads
    \item \textbf{One-Side vs Both-Side Detection:} Ratio-based method has ambiguous zone (0.8-1.2) where classification is uncertain
    \item \textbf{Sidewalk Gaps:} Small gaps in sidewalk network (driveways, etc.) may cause discontinuity in coverage calculations
    \item \textbf{Intersection Handling:} Complex intersections may cause double-counting or misclassification
\end{enumerate}

% ============================================================================
\section{Output Data Products}
% ============================================================================

\subsection{Generated Datasets}

All analysis methods produce the following standardized outputs:

\subsubsection{GeoJSON Files}

\begin{itemize}
    \item \textbf{Complete Coverage:} All roads with coverage classification
    \item \textbf{By Coverage Type:} Separate files for none/one\_side/both\_sides
    \item \textbf{Coordinate System:} EPSG:4326 (WGS84) for web mapping compatibility
\end{itemize}

\subsubsection{CSV Files}

Tabular export of results (geometry dropped) for spreadsheet analysis and database import.

\subsubsection{JSON Statistics}

Comprehensive statistics including:
\begin{itemize}
    \item Coverage counts and percentages
    \item By-road-type breakdowns
    \item Method-specific metrics (e.g., buffer distances)
\end{itemize}

\subsection{Output Directory Structure}

\begin{verbatim}
Technical/data/processed/sidewalk_method_comparison/
  method1_dvrpc_baseline/
    dvrpc_coverage.geojson
    roads_coverage_none.geojson
    roads_coverage_one_side.geojson
    roads_coverage_both_sides.geojson
    dvrpc_coverage.csv
    dvrpc_statistics.json
  method2_adaptive_buffer/
    adaptive_buffer_coverage.geojson
    ...
  method3_edge_detection/
  method4_centerline_perpendicular/
  method5_network_voronoi/
\end{verbatim}

% ============================================================================
\section{Recommendations for Future Enhancement}
% ============================================================================

\subsection{Data Collection}

\begin{enumerate}
    \item \textbf{Complete Road Centerline Dataset:} Obtain or generate complete centerline network matching all 4,386 road polygons
    \item \textbf{Sidewalk Width Attributes:} Collect sidewalk width data for ADA compliance analysis
    \item \textbf{Crosswalk and Curb Ramp Data:} Add intersection accessibility data
    \item \textbf{Maintenance Records:} Integrate sidewalk condition and repair history
\end{enumerate}

\subsection{Analysis Enhancements}

\begin{enumerate}
    \item \textbf{True Centerline Extraction:} Implement medial axis/skeletonization algorithm for accurate centerlines
    \item \textbf{Variable Width Modeling:} Develop method to handle roads with significant width variation along length
    \item \textbf{Improved Side Detection:} Implement left/right classification using perpendicular transects or Voronoi diagrams
    \item \textbf{Network Analysis:} Add connectivity analysis for pedestrian routing
    \item \textbf{Equity Analysis:} Integrate demographic data for environmental justice assessment
\end{enumerate}

\subsection{Validation}

\begin{enumerate}
    \item \textbf{Ground Truth Sample:} Field verify classification for random sample of roads
    \item \textbf{Aerial Imagery Validation:} Visual comparison with high-resolution imagery
    \item \textbf{Method Comparison:} Implement multiple methods and compare agreement
    \item \textbf{Sensitivity Analysis:} Test impact of buffer distance and threshold parameters
\end{enumerate}

% ============================================================================
\section{Conclusion}
% ============================================================================

This analysis utilizes high-quality geospatial data from Westchester County to comprehensively assess sidewalk coverage across the county's 4,386-segment road network. All input data is well-documented, with clear coordinate systems, consistent geometries, and comprehensive spatial coverage.

Key transformations include:
\begin{itemize}
    \item Algorithmic centerline extraction from road polygons
    \item Adaptive buffer creation based on calculated road width
    \item Spatial indexing for efficient intersection detection
    \item Ratio-based coverage classification (none/one-side/both-sides)
\end{itemize}

The analysis makes full use of geometric data while identifying numerous attribute fields and external datasets that could enhance future iterations. No data is filtered or excluded, ensuring comprehensive county-wide assessment.

% ============================================================================
\section{Appendix: File Inventory}
% ============================================================================

\subsection{Input Files}

\begin{longtable}{p{5cm}p{2cm}p{6cm}}
\toprule
\textbf{Filename} & \textbf{Size} & \textbf{Description} \\
\midrule
\endhead
countyroadways\_polygon.shp & 194.0 MB & Road polygon geometries \\
countyroadways\_polygon.dbf & 469.5 KB & Road polygon attributes \\
countyroadways\_polygon.shx & 35.2 KB & Road polygon spatial index \\
countyroadways\_polygon.prj & 482 B & Projection file (EPSG:2260) \\
countyroadways\_polygon.xml & 47.5 KB & Metadata \\
countyroads\_line.shp & 16.1 KB & Road centerline geometries \\
countyroads\_line.dbf & 9.4 KB & Road centerline attributes \\
countyroads\_line.shx & 468 B & Road centerline spatial index \\
countyroads\_line.prj & 482 B & Projection file (EPSG:2260) \\
countyroads\_line.xml & 52.1 KB & Metadata \\
countysidewalks\_polygon.shp & 87.3 MB & Sidewalk polygon geometries \\
countysidewalks\_polygon.dbf & 8.3 MB & Sidewalk polygon attributes \\
countysidewalks\_polygon.shx & 620.6 KB & Sidewalk polygon spatial index \\
countysidewalks\_polygon.prj & 482 B & Projection file (EPSG:2260) \\
countysidewalks\_polygon.xml & 47.5 KB & Metadata \\
\midrule
\textbf{Total} & \textbf{282.0 MB} & \textbf{15 files} \\
\bottomrule
\caption{Complete Input File Inventory}
\label{tab:file_inventory}
\end{longtable}

\subsection{Processing Scripts}

\begin{longtable}{p{6cm}p{7cm}}
\toprule
\textbf{Script} & \textbf{Purpose} \\
\midrule
\endhead
method1\_dvrpc\_baseline.py & DVRPC fixed-buffer baseline analysis \\
method2\_adaptive\_buffer.py & Adaptive buffer based on road width \\
method3\_edge\_detection.py & Edge-based left/right detection \\
method4\_centerline\_perpendicular.py & Perpendicular transect method \\
method5\_network\_voronoi.py & Network Voronoi diagram approach \\
sidewalk\_method\_comparison.py & Multi-method comparison and validation \\
\bottomrule
\caption{Analysis Scripts Inventory}
\label{tab:scripts}
\end{longtable}

% ============================================================================
% Bibliography
% ============================================================================

\begin{thebibliography}{99}

\bibitem{dvrpc2017}
Delaware Valley Regional Planning Commission (DVRPC),
\textit{Pedestrian Facilities Network: Sidewalk Inventory Methodology},
2017.

\bibitem{nacto2013}
National Association of City Transportation Officials (NACTO),
\textit{Urban Street Design Guide: Sidewalk Design Standards},
2013.

\bibitem{fhwa2024}
Federal Highway Administration (FHWA),
\textit{Guide for Scalable Risk Assessment Methods for Pedestrians and Bicyclists},
2024.

\bibitem{esri2025}
ESRI,
\textit{ArcGIS Pro Documentation: Buffer Analysis Best Practices},
2025.

\bibitem{shapely}
Shapely Development Team,
\textit{Shapely Documentation: Geometric Operations},
\url{https://shapely.readthedocs.io/},
2025.

\end{thebibliography}

\end{document}
